\section{SycoScope: Sycophancy Probe using Mechanistic Interpretability}\label{sec:sycoscope}% Draft Oscar:
\textbf{Linear Probing}
We train distinct linear probes on both the hidden states of the residual stream and the output of attention for each trigger type $t$ and layer $\ell$. This produces $t=8$ distinct linear probes $$f_t = \sigma(w_tx+b_t)$$ For each linear probe, we calculate and compare the average accuracy and AUC-ROC in a holdout validation set to determine the best location and layer to utilize weights for subspace analysis. 

\textbf{Sycophancy Subspace}
\cite{shahgeometry2025} demonstrated that different harmful subjects are encoded in a subspace of low-rank in the attention output. Furthermore, \cite{vennemeyer2025causalsycophancy} showed that sycophantic agreement and sycophantic praise occupy separable subspaces that can be amplified or suppressed independently of each other, establishing that distinct subtypes of sycophancy have geometrically distinct representations. We ask whether sycophancy, as a distinct failure mode, exhibits an analogous structure across trigger types. As an initial exploratory analysis, we first normalize the weights for each trigger probe $w^*_t = \frac{w_t}{\vert \vert w_t \vert \vert}$ and stack each weight to $W \in \mathbb{R}^{t \times d}$. We then measure the linear dimensionality of the stacked normalized weight matrix $W$ using the effective rank $k$. Given an energy threshold $\tau$, we calculate $k$ as the minimum number of orthogonal components that is needed to explain $\tau$ percent of the variance in the weight matrix. $$k = \min \{r: \frac{\sum_{i = 1}^r\sigma^2_i}{\sum_{i = 1}^n \sigma^2_i}\geq \tau\}$$ A low effective rank indicates that different trigger types share a common subspace in the residual stream. A high effective rank indicates that each trigger type occupies a distinct direction.

\textbf{Causal Validation}
To determine whether the identified geometry causally encodes sycophancy, we first steer using the probe weight vectors $w^*_t$ directly. If steering along $w^*_t$ induces sycophantic behavior under a neutral prompt, probe weights causally encode sycophancy and we proceed using $W$. Otherwise, we construct difference-in-means trigger contrast vectors $$c_t = \frac{1}{|\mathcal{D}_t^+|}\sum_{x_i^+} \mathbf{h}^{\ell^*}(x_i^+) - \frac{1}{|\mathcal{D}_t^-|}\sum_{x_j^-} \mathbf{h}^{\ell^*}(x_j^-)$$ where $\mathcal{D}_t^+$ and $\mathcal{D}_t^-$ are the sets of sycophantic and non-sycophantic responses conditioned on trigger type $t$, and $\mathbf{h}^{\ell^*}(x) \in \mathbb{R}^d$ is the activation of the residual stream at layer $\ell^*$. We causally test each $c_t$ independently via steering, then stack the normalized vectors into $W_{\text{DIM}} \in \mathbb{R}^{t \times d}$ and recompute the effective rank. Subsequent analysis proceeds on the matrix that causally encodes sycophancy.

\textbf{High Rank Results}
To visualize the geometric separation between trigger types, we project the residual stream activations at $\ell^*$ onto their top two principal components via PCA and also apply t-SNE to $d$-dimensional activations. Distinct clusters in these projections corroborate the high-rank finding, while overlap would suggest that the effective rank overestimates true geometric separation.

\textbf{Intermediate Rank Results}
To test whether the geometries are similar, we computed a pairwise cosine similarity matrix between each vector. In the intermediate case, the similarity matrix reveals a clustering structure to test which trigger families share geometric spaces that effective rank along collapses. 

\textbf{Low Rank Results}
If $W$ is approximately low-rank, different trigger types share a common sycophancy subspace. We steer along the top singular vector $v_1 \in \mathbb{R}^d$ of $W$ under a neutral prompt. Successful elicitation of sycophantic behavior causally validates that the shared subspace is behaviorally sufficient, rather than an artifact of probe training. 

\textbf{TunedLens for Layer-wise Evolution} Previous work suggests that transformer predictions follow a two-stage structure: an early phase in which the model's layer-wise logit distribution shifts toward a truthful prediction, followed by a late phase in which sycophancy overrides truth \cite{wang2025truthoverridden}. To locate this override, we apply tuned lens and logit lens to track the layer-wise evolution of the sycophantic response probability across the residual stream, asking whether sycophantic commitment emerges earlier under high-severity triggers or under particular trigger families.

\textbf{Causal Validation of Capitulation Layer}
To causally validate the two-stage override, we perform two complementary last-token patching experiments. First, we patch residual stream activations at the last token from a triggered clean run that elicits a truthful response into a corrupted run that elicits sycophancy, sweeping from early to late layers to identify the critical layer $\ell^*$ beyond which the model can no longer recover a truthful answer for each trigger family. Since both runs contain the trigger but differ only in whether the model capitulates, this design isolates the sycophantic commitment signal independent of trigger encoding. Second, we patch last-token activations from a neutral prompt run into a corrupted triggered run, sweeping across layers to identify the earliest layer at which the trigger's representational influence becomes causally sufficient to produce sycophancy. Together, these two designs disentangle trigger encoding from behavioral commitment, providing a layerwise causal account of where sycophancy first becomes detectable and where it becomes irreversible.

